% Generated draft: requires mathematical and source review.
\documentclass[11pt]{article}
\usepackage[margin=25mm]{geometry}
\usepackage{fontspec}
\setmainfont{Latin Modern Roman}
\usepackage{amsmath,amssymb,mathtools,unicode-math}
\setmathfont{Latin Modern Math}
\usepackage{enumitem,needspace,xurl,hyperref}
\setlength{\parindent}{0pt}
\setlength{\parskip}{5pt}
\setlength{\emergencystretch}{3em}
\newcommand{\N}{\mathbb{N}}
\newcommand{\Z}{\mathbb{Z}}
\newcommand{\Q}{\mathbb{Q}}
\newcommand{\R}{\mathbb{R}}
\newcommand{\C}{\mathbb{C}}
\newcommand{\F}{\mathbb{F}}
\newcommand{\A}{\mathbb{A}}
\newcommand{\PP}{\mathbb{P}}
\newcommand{\E}{\mathbb{E}}
\newcommand{\eps}{\varepsilon}
\newcommand{\dd}{\,\mathrm{d}}
\begin{document}
\section*{UnsolvedMath ID: 20002723}
\textbf{Source problem code:} catalogue entry
\subsection{Definitions and mathematical statement}
A. Zelevinsky's question asks for an exceptional-type visual interpretation of $\Delta_W$ analogous to the polygon/triangulation models in classical types $A,B,C,D$. Interpret $\Delta_W$ as the cluster complex: vertices are almost-positive roots and facets are clusters. Calling it an associahedron can refer to a dual polytope convention, which must be distinguished. The source does not require the exceptional model to be one ordinary undecorated polygon.

The aid reports substantial later colored-diagonal and cluster-category interpretations, most completely for $E_6$ and its $F_4$ folding, requiring colors, orientations, paired objects, several polygon layers, and sometimes mutations that are not ordinary flips. It supplies no complete model or bibliographic references here.

Its first result is a counting obstruction to the most restrictive polygon interpretation. For an irreducible rank-$n$ finite root system with Coxeter number $h$, $\Delta_W$ has $n(h+2)/2$ vertices and facets of size $n$. The ordinary diagonal complex of a convex $N$-gon has $N(N-3)/2$ vertices and triangulation facets of size $N-3$. A simplicial isomorphism therefore requires
\[
 N=n+3,\qquad \frac{n(h+2)}2=\frac{n(n+3)}2,
 \quad\text{hence }h=n+1\quad(n>0).
\]
This excludes $G_2,F_4,E_6,E_7,E_8$ from being the triangulation complex of a single undecorated polygon. It does not exclude decorated, folded, layered, or other geometric models, and the numerical condition alone is not a sufficient isomorphism criterion.

The second result is an abstract finite folding certificate. Let a finite group act simplicially on a finite flag complex $K$. A vertex orbit $O$ is admissible if its vertices form a face. Two admissible orbits are compatible if their union is a face. Define the folded complex on admissible orbits by requiring their total union to be a face. Because $K$ is flag, internal admissibility and pairwise cross-orbit compatibility exactly determine its faces. Its faces correspond to group-invariant faces of $K$ by taking their unions.

For facets, an extra symmetric-extension property is required: every invariant face that is not already a facet of $K$ can be enlarged by an entire admissible orbit while remaining a face. In the finite setting this is equivalent to every invariant face extending to an invariant facet, and exactly ensures that folded facets lift to invariant facets rather than merely maximal invariant nonfacets. A concrete exceptional-model certification must also identify the relevant root/orbit compatibility, not merely match the number of facets. The aid does not provide those identifications for all exceptional types, so the general visual-model request remains broader than these two audit criteria.
\subsection{Short English statement}
An ordinary undecorated polygon cannot realize the listed exceptional cluster complexes because facet and vertex counts force the type-$A$ Coxeter-number relation. The supplied folding criterion explains what extra compatibility and extension checks a decorated model needs, without constructing all exceptional models.
\subsection{Sources}
\begin{itemize}
\item Local UnsolvedMath registry, permanent ID 20002723, source code catalogue entry; created and updated 14 August 2026. Canonical \texttt{aim-probability-notes.json}, zero-based \texttt{notes[164]}, tag \texttt{problem}.
\item AIM ``Braid groups, clusters and free probability,'' January 2005 report, section 6, Problem 6.2 attributed to A. Zelevinsky, physical page 12 according to the supplied verification: \url{https://aimath.org/WWN/braidgroups/braidgroups.pdf}. The clean field restores the split heading; statement status is reconstructed\_unverified. Original remarks and literature field are empty.
\item Research attempt 1: \texttt{partial\_result}, problem status at run \texttt{partially\_solved}, rigor label \texttt{proved}, novelty confidence low. The later colored-diagonal/cluster-category constructions and the obstruction/folding results are not accompanied by complete authoritative citations or detailed model artifacts.
\item Ulam AI / UnsolvedMath Contributors, CC BY 4.0. Machine-generated research aid, not peer reviewed; independently verify before citation.
\end{itemize}
\end{document}
